% !TEX root = ../Vorlage_DA.tex

%	########################################################
% 			INFO: Zitieren, Abbildungen, Listing
%	########################################################


%	--------------------------------------------------------
% 	Überschrift, Inhaltsverzeichnis
%	--------------------------------------------------------
\chapter{INFO: Zitieren, Abbildungen, Quelltext\authorA}\label{ref:zitieren}

In diesem Kapitel sind Beispiele angeführt, wie Abbildungen, Zitate und Quelltext zu verwenden sind.

Die Zitate, Abbildungen und Listings werden automatisch in das Quellen- und Abbildungsverzeichnis übernommen. 
Das Literatur-, Abbildungs- und Listingsverzeichnis sind am Ende der Arbeit zu finden.   

%	--------------------------------------------------------
% 	Inhalt 
%	--------------------------------------------------------
\section{Abbildungen}\label{ref:abbildungen}

Abbildung sind mit einer Abbildungsnummer und einer Unterschrift zu versehen die kurz die Abbildung beschreibt. 
Abbildungen gehören zum umgebenden Text und müssen dort erwähnt werden.

Beispiel: In Abbildung \ref{htl01} wird das Logo der HTL Braunau dargestellt.

Wird auf eine Abbildung referenziert die sich weiter von der aktuellen Textstelle entfernt befindet so kann auch die Seitennummer hinzugefügt werden.

Beispiel: Siehe Abbildung \ref{htl01} auf Seite \pageref{htl01}.


%	--------------
%	Eine Abbildung
\begin{figure}[H]
	\centering
	\includegraphics[width=0.5\textwidth]{./media/images/htl_c_cmyk_rein.pdf}
  	\caption{Logo der HTL Braunau.}
  	\label{htl01}
\end{figure}


%	----------------
%	Zwei Abbildungen

\begin{figure}[H]
  \centering
  \subfloat[Logo der HTL Braunau.]{\label{fig:a}\includegraphics[width=.25\textwidth]{./media/images/htl_c_cmyk_rein.pdf}}
  \qquad
  \subfloat[Logo der HTL Braunau.]{\label{fig:b}\includegraphics[width=.25\textwidth]{./media/images/htl_c_cmyk_rein.pdf}}
  \label{fig:canvas_01}
  % caption - In eckigen Klammern: Text für das Abbildungsverzeichnis
  \caption[Zwei Logos der HTL Braunau]{Zwei Logos der HTL Braunau in einer Abbildung zusammengefasst.}
\end{figure}


%	----------------
%	Drei Abbildungen
\begin{figure}[H]
  \centering
  \subfloat[Logo]{\label{fig:a}\includegraphics[width=.15\textwidth]{./media/images/htl_c_cmyk_rein.pdf}}
  \qquad
  \subfloat[Logo]{\label{fig:b}\includegraphics[width=.15\textwidth]{./media/images/htl_c_cmyk_rein.pdf}}
\qquad
  \subfloat[Logo]{\label{fig:b}\includegraphics[width=.15\textwidth]{./media/images/htl_c_cmyk_rein.pdf}}
  \label{fig:canvas_04}
\caption[Drei Logos der HTL Braunau]{Dreimal das Logo der HTL Braunau.}
\end{figure}


\section{Zitate, Quellen, Fußnoten}

Es muss dem Leser möglich sein alle dargestellten Informationen selbst zu überprüfen.
Daher gilt der wissenschaftliche Grundsatz: 
Wer Informationen/Erkenntnisse verwendet die nicht von einem selbst stammen muss dies eindeutig kennzeichnen.
Das gilt für Bücher, Zeitschriftenartikel aber auch für alle Internetquellen.

Dem Leser Hinweise darüber zu geben wo genauere/weitere Informationen zu einem Thema zu finden sind ist ein weiterer Grund für Zitieren.

Es ist erlaubt fremde Quellen zu zitieren solange dies eindeutig erkennbar ist und sich die zitierten Stellen auf einen kurzen Absatz mit wenigen Zeilen beschränken. Auch wenn nicht wortwörtlich zitiert wird ist die Quelle der Information anzugeben.

Wörtliche Zitate werden durch Anführungszeichen begonnen und beendet sowie kursiv geschrieben: 
\emph{\glqq Wissenschaftliches Plagiat: Man kann sich zwar mit fremden Federn schmücken, aber man kann nicht mit ihnen fliegen.\grqq}~\cite{bib:uhlenbruck1}


Längere Zitate werden durch Einrücken vom normalen Text abgesetzt:
\begin{quotation}
\noindent
\emph{\glqq We tend to think of navigating a website as clicking from page-to-page via some kind of global navigation that’s always visible. When it comes to a single page, we often think scrolling is the one and only way to move from one end to the next.\grqq}~\cite{bib:bradley1}
\end{quotation}

Die eckigen Klammer mit Nummer kennzeichnet die Quelle.
Am Ende des Dokuments werden alle Quellen im Literaturverzeichnis aufgelistet. \cite{bib:wikilitverz}.

Das Verwenden fremden Gedankenguts ohne die Quelle anzugeben ist ein Plagiat (geistiger Diebstahl, \cite{bib:wikiplagiat}) und unter Umständen sogar eine Urheberrechtsverletzung.


Verweise in das Literaturverzeichnis sind nicht auf Zitate beschränkt sondern können auch eingesetzt werden um darauf hinzuweisen wo zu einem Thema mehr Informationen erhältlich sind.


Es ist zunehmend eine Kurzzitierweise in Fußnoten üblich: Nachname des Verfassers, Kurztitel, Seitenangabe.\footnote{John Doe, \glqq Verwendung von Fußnoten\grqq, Seite 12.}


%	%%%%%%%%%%%%%%%%%%%%%%%%%%%%%%%%%%%%%%%%%%%%%
%	Listings, Code
%	%%%%%%%%%%%%%%%%%%%%%%%%%%%%%%%%%%%%%%%%%%%%%
\section{Listings, Code}

Quelltext ist ein wesentlicher Bestandteil einer Informatik Diplomarbeit.
Listings werden wie Abbildungen nummeriert und mit einer Unterschrift versehen.
Ebenfalls müssen sie im Text referenziert werden.

\subsection{Beispiele}

Listing \ref{code:helloworld} auf Seite \pageref{code:helloworld} zeigt ein Hallo Welt Programm.

%	%%%%%%%%%%%%%%%%%%%%%%%%%%%%%%%%%%%%%%%%%%%%%
%	C++
\begin{minipage}{\linewidth}
\lstinputlisting[language=C++,caption={Mein erstes C-Programm.},label=code:helloworld]{./media/code_snippets/bsp_01.cpp}
\end{minipage}


Das PHP Programm \ref{code:phpread} dient zum ermitteln aller Dateien eines Unterverzeichnisses.

\begin{lstlisting}[language=PHP,caption={Alle Dateinamen ausgeben.},label=code:phpread]
<?php
if ($handle = opendir(realpath(__DIR__ . '/../userimages'))) {
    
    while (false !== ($file = readdir($handle))) {
        echo "$file//";
    }

    closedir($handle);
}
?>
\end{lstlisting}

\begin{lstlisting}[language=JavaScript,caption={Mausposition ermitteln.}]
// Calculates the mouse position.
function getMousePos(canvas, evt) {
	var rect = canvas.getBoundingClientRect();
	return {
		x: evt.clientX - rect.left,
		y: evt.clientY - rect.top
	};
}
\end{lstlisting}




%	%%%%%%%%%%%%%%%%%%%%%%%%%%%%%%%%%%%%%%%%%%%%%
%	CSS
\lstinputlisting[language=CSS,caption={Kurzer CSS--Quelltext.}]{./media/code_snippets/bsp.css}


%	%%%%%%%%%%%%%%%%%%%%%%%%%%%%%%%%%%%%%%%%%%%%%
%	HTML
\lstinputlisting[language=HTML,caption={HTML--Quelltext}]{./media/code_snippets/bsp.html}


%	%%%%%%%%%%%%%%%%%%%%%%%%%%%%%%%%%%%%%%%%%%%%%
%	HTML
\lstinputlisting[language=JavaScript,caption={JavaScript--Quelltext}]{./media/code_snippets/bsp.js}



